% ---------------------------------------------------------------------------
%  Plantilla de presentacion -- Universidad de Santiago de Chile
%  Tema: usach v2.0    Compilar: latexmk -pdf -outdir=build main.tex
%  (xelatex/lualatex activan las fuentes corporativas si estan instaladas:
%   latexmk -xelatex -outdir=build main.tex)
% ---------------------------------------------------------------------------
\documentclass[11pt,aspectratio=169]{beamer}

% --- Idioma: babel-spanish si esta disponible; si no, ajustes manuales -------
\IfFileExists{spanish.ldf}{%
  \usepackage[spanish,es-noquoting,es-nodecimaldot]{babel}%
}{%
  \usepackage[T1]{fontenc}
  \renewcommand{\figurename}{Figura}
  \renewcommand{\tablename}{Tabla}
  \renewcommand{\refname}{Referencias}
}
\usepackage[utf8]{inputenc}   % inofensivo en XeLaTeX/LuaLaTeX modernos
\usepackage{amsmath,amssymb}
\usepackage{booktabs}
\usepackage{microtype}

\usetheme{usach}              % opciones: sinprogreso, logolamina, sinportadillas, pei

% --- Imagotipo oficial ------------------------------------------------------
% Ya copiados a assets/ desde el ZIP oficial (guiaweb.usach.cl).
%   \usachimagotipo{assets/Usach S1.png}       % positivo (fondo claro)
%   \usachimagotipoblanco{assets/Usach SB.png} % negativo (fondo oscuro)
% Alternativa vertical (prioritaria MNG): assets/Usach P2.png y Usach PB.png.
\usachimagotipo{assets/Usach S1.png}
\usachimagotipoblanco{assets/Usach SB.png}

% --- Metadatos --------------------------------------------------------------
\title[Anomalías en telemetría]{Detección de anomalías en telemetría de medidores}
\subtitle{Un modelo de series de tiempo para lecturas de consumo en terreno}
\author{Iván Ahumada}
\date{Septiembre de 2026}

\usachfacultad{Facultad de Ingeniería}
\usachdepartamento{Departamento de Ingeniería Informática}
\usachprograma{Magíster en Ingeniería Informática}
\usachasignatura{Computación Científica}
\usachprofesor{Prof. Pablo Román}

% Texto corto del pie (izquierda). Si se omite, se usa el titulo completo.
\titlegraphic{}
\institute[USACH]{Universidad de Santiago de Chile}

\begin{document}

\usachportada

% --- Agenda ------------------------------------------------------------------
\begin{frame}{Contenido}
  \tableofcontents
\end{frame}

\section{Problema}

\begin{frame}{Qué falla hoy}{Lecturas de terreno sin control de calidad automático}
  \begin{itemize}\usachaire
    \item El 8\,\% de las lecturas mensuales llega fuera del rango físico esperable.
    \item La validación es manual y ocurre después de emitir la boleta.
    \item Cada corrección posterior cuesta una revisita a terreno.
  \end{itemize}

  \vfill
  \usachdatos{%
    \begin{minipage}[t]{0.31\linewidth}\usachdato{8\,\%}{lecturas fuera de rango}\end{minipage}\hfill
    \begin{minipage}[t]{0.31\linewidth}\usachdato{14 días}{demora media de detección}\end{minipage}\hfill
    \begin{minipage}[t]{0.31\linewidth}\usachdato{1,7}{revisitas por corrección}\end{minipage}}
\end{frame}

\begin{frame}{Bloques del sistema}
  \begin{columns}[T,onlytextwidth]
    \begin{column}{0.48\textwidth}
      \begin{block}{Definición}
        Una lectura es anómala si se aparta del consumo esperado del predio más
        allá de tres desviaciones robustas.
      \end{block}
      \begin{exampleblock}{Resultado esperado}
        Marcar la lectura antes del cierre del ciclo de facturación.
      \end{exampleblock}
    \end{column}
    \begin{column}{0.48\textwidth}
      \begin{alertblock}{Restricción}
        El modelo corre en el dispositivo del lector, sin conexión garantizada.
      \end{alertblock}
      \usachcita{La calidad del dato se define en el punto de captura, no en el
      escritorio.}{Nota de terreno, cuadrilla norte}
    \end{column}
  \end{columns}
\end{frame}

\section{Método}

\usachdestacado{Si el dato se corrige en terreno, el ciclo de facturación no
se vuelve a abrir.}

\begin{frame}{Formulación}
  Sea $y_t$ la lectura acumulada en el período $t$ y $\Delta y_t = y_t - y_{t-1}$
  el consumo del período. El residuo robusto es
  \[
    r_t \;=\; \frac{\Delta y_t - \operatorname{med}(\Delta y)}{1{,}4826 \cdot \operatorname{MAD}(\Delta y)},
  \]
  y se marca la lectura cuando $|r_t| > 3$.

  \begin{itemize}
    \item La mediana y el MAD se recalculan con las últimas doce lecturas del predio.
    \item El umbral es configurable por tipo de cliente.
  \end{itemize}
\end{frame}

\begin{frame}{Comparación de estrategias}
  \begin{table}
    \centering
    \small
    \begin{tabular}{@{}lccc@{}}
      \toprule
      Estrategia & Precisión & Recall & Latencia \\
      \midrule
      Rango fijo por tarifa      & 0,41 & 0,88 & inmediata \\
      Residuo robusto (MAD)      & 0,79 & 0,84 & inmediata \\
      Modelo estacional ajustado & 0,86 & 0,81 & 2 s \\
      \bottomrule
    \end{tabular}
    \caption{Desempeño sobre 24 meses de lecturas etiquetadas.}
  \end{table}
  \usachfuente{elaboración propia a partir de la base de lecturas 2024--2026.}
\end{frame}

\section{Resultados}

\begin{frame}{Lo que cambió}
  \begin{itemize}\usachaire
    \item La detección pasó de 14 días a la misma visita.
    \item Las revisitas por corrección bajaron a 0,4 por lectura marcada.
    \item El modelo cabe en 40\,kB y corre sin conexión.
  \end{itemize}
  \vfill
  \begin{block}{Próximo paso}
    Validar el umbral por tipo de cliente en la zona sur antes de extenderlo.
  \end{block}
\end{frame}

\begin{frame}[allowframebreaks]{Referencias}
  \begin{thebibliography}{9}
    \bibitem{rousseeuw1993}
      Rousseeuw, P. y Croux, C. (1993).
      \newblock Alternatives to the median absolute deviation.
      \newblock \emph{Journal of the American Statistical Association}, 88(424).
    \bibitem{mng2024}
      Universidad de Santiago de Chile (2024).
      \newblock \emph{Manual de Normas Gráficas 2024-2S}.
  \end{thebibliography}
\end{frame}

\usachcierre{ivan.ahumada@usach.cl \\ Departamento de Ingeniería Informática}

\appendix

\begin{frame}{Anexo: parámetros del modelo}
  \begin{itemize}
    \item Ventana móvil: 12 períodos.
    \item Factor de consistencia: 1,4826 (normalidad asintótica).
    \item Umbral por defecto: $|r_t| > 3$.
  \end{itemize}
\end{frame}

\end{document}
